
\blankpage
\pagebreak

{
\thispagestyle{empty}

\mbox{}

\vspace{11cm}

\footnotesize
\noindent{}\emph{\noindent{}Conta Martin Buber: ``Perguntaram ao Rabi Pinkhas: Por que é que, como afirma a tradição, o Messias deverá nascer no dia do aniversário da destruição do Templo? O grão semeado na terra --- replicou ele --- deve apodrecer, para que brote a nova espiga. O poder não pode ressuscitar, sem antes dissolver-se no grande segredo. Despir uma forma, assumir uma forma, isso se faz no instante do puro nada. Na concha do esquecimento, cresce o poder da memória. No dia da destruição, o poder jaz no fundo e aí cresce. É por isso que nesse dia nos sentamos no chão. É por isso que nesse dia visitamos os túmulos. É por isso que nesse dia há de nascer o Messias.''\footnote{\textsc{buber}, \emph{op.\,cit.}, p.166.}}
}

\pagebreak

\chapter{Lugares}

% \noindent{}Conta Martin Buber: ``Perguntaram ao Rabi Pinkhas: Por que é que,
% como afirma a tradição, o Messias deverá nascer no dia do aniversário da
% destruição do Templo? O grão semeado na terra --- replicou ele --- deve
% apodrecer, para que brote a nova espiga. O poder não pode ressuscitar,
% sem antes dissolver-se no grande segredo. Despir uma forma, assumir uma
% forma, isso se faz no instante do puro nada. Na concha do esquecimento,
% cresce o poder da memória. No dia da destruição, o poder jaz no fundo e
% aí cresce. É por isso que nesse dia nos sentamos no chão. É por isso que
% nesse dia visitamos os túmulos. É por isso que nesse dia há de nascer o
% Messias''\footnote{\textsc{buber}, \emph{op.\,cit.}, p.166.}.

Os casamentos judaicos podem ser realizados em quaisquer lugares,
contanto que sejam feitos debaixo de uma \emph{chupá}: um pedaço de pano
retangular ou quadrado tensionado, formando uma espécie de cabana, e
sustentado por quatro pilares, cada qual em uma extremidade. A
\emph{chupá} é um espaço sagrado em que o casal se encontra para selar a
união entre si e com Deus. Finaliza-se este ritual com a quebra de um
copo de vidro sob o chão pisado pelo par. Simultaneamente, ou seja, no
exato instante desta ruptura, o pedaço de pano --- que representa o
lugar santo da \emph{chupá} --- é recolhido. Já não há mais espaço
sagrado, o copo foi quebrado e a \emph{Shechiná} --- a habitação e a
presença de Deus que estava dentro do território temporário e apropriado
da \emph{chupá} --- foi, em forma de vento e em forma de sopro,
espalhada pelo universo.

O espraiamento e a dispersão\footnote{A palavra ``dispersão'' provém da
  raiz grega ``diasperein'': dispersar, espalhar, disseminar. O prefixo
  ``dia'' (através de) unido à raiz ``sperein'' (esporos ou sementes)
  revela, ao final, que ``diasperein'' significa através de sementes ou
  de esporos. Dessa forma, toda dispersão é uma semeadura, ou seja, tudo
  nasce para que seja espalhado/\,dispersado e levado pelo velho vento do
  sopro nômade da diáspora para que semeie e que brote em outras
  localidades.} da \emph{Shechiná} pelo mundo se apresentam enquanto um
momento de uma lamentação tão profunda que se transforma em festividade.
Assim como é, também, a música iídiche. O estilo
\emph{klezmer}\footnote{\emph{Klezmer} é uma palavra de origem
  iídiche que significa instrumentos musicais, e é também um
  gênero musical não litúrgico desenvolvido em meados do século \textsc{xv} pela
  cultura \emph{ashkenazi} (judeus provenientes da Europa Central e do
  Leste Europeu), em que a dança é um elemento determinante e as
  melodias imitam os risos, choros e suspiros dos \emph{hassidim}. \emph{Hassidim} (plural de \emph{hassid}) são os membros do
  \emph{hassidismo} judaico, um movimento que dá ênfase à alegria
  religiosa e é organizado por dinastias, tendo sempre um líder
  espiritual em cada uma das gerações.} são melodias, cantos e cantigas
de tanta melancolia e lamentação que se transformam em espaços de
celebração de vida.

A quebra do copo de vidro é uma recordação de que judeus e judias estão
incompletos sem a presença do Templo Sagrado de Jerusalém. Porém,
deve-se ter em mente que todos os eventos evocadores do Templo se firmam
enquanto momentos de celebração, assim como é o casamento.

Anteriormente à quebra do copo lê-se o Salmo 137, que diz: 

\begin{verse}
Se eu me esquecer de Ti, ó Jerusalém\\
Que minha mão direita resseque\\
Que minha língua cole no céu da boca\\
Se eu não lembrar de Ti\\
Se eu não fizer de Ti, ó Jerusalém\\
a maior das minhas alegrias.
\end{verse}

Este versículo é uma afirmação de devoção e de vida das comunidades judaicas, onde
constantemente se rememora a incompletude com a ausência do Templo
Sagrado, mesmo nos momentos de maior felicidade. Assim, como diz Rabi
Pinkhas: ``{[}\ldots{}{]} No dia da destruição o poder jaz no fundo e
aí cresce''\footnote{\textsc{buber}, \emph{op.\,cit.}, p.166.}. É nestes momentos de
ruptura, de quebra, de lamentação e de melancolia que os povos em
diáspora, antes de qualquer coisa, celebram a vida.

Comentam os sábios que a santidade espiritual atingida no momento da
\emph{chupá} é tamanha que se encontram neste local sagrado --- e
temporário --- as dez últimas gerações de pessoas falecidas que fazem
parte destes fios, destes nós e destas linhas genealógicas. As almas
descem (e descendem\footnote{Se partirmos do pressuposto de que tudo vem
  do pó, para retornar ao pó, então as almas retornam, descem, e ao
  mesmo tempo descendem, da esfera divina. A vida e a morte estão no
  mesmo lugar.}) da esfera divina e agregam-se no local da \emph{chupá}
abençoando e santificando a união do casal.

Por fim, a \emph{chupá} cria um lugar santo e temporário, transitório e
apropriado em que se encontram, ao mesmo tempo, Deus e toda a sua
criação. Deus é toda a criação (\emph{HaMakom}\footnote{\emph{HaMakom}
  significa, em português, ``o Lugar'' e é um dos nomes judaicos de
  Deus.}), é o lugar por completo que engloba os vivos e os mortos. Deus
envolve não somente os lugares em que estamos --- em que vivemos ---,
mas toda a existência em si. E \emph{HaMakom} envolve tudo, assim como o
útero envolve o feto.

Durante a santidade do momento do \emph{shabat} clama-se: ``Vem,
meu Amado, saudar a Noiva, vem, vamos dar as boas-vindas à Presença do
shabat''. A canção de \emph{Lechá Dodi}\footnote{O \emph{Lechá Dodi} é
  um poema cantado e entoado no momento de início do \emph{shabat}. Esta
  cantiga é de autoria do Rabi Shlomo Halevi Alkabets, composta entre os
  anos de 1529 e 1540. Porém foi Isaac Luria (1534--1572), conhecido como
  Ari, que desempenhou um papel fundamental para a popularização e
  espiritualização do hino.} é o ponto máximo da elevação espiritual do
\emph{cabalat shabat} (literalmente ``o recebimento do \emph{shabat}''), realizado às
sextas-feiras à noite. Nesse momento, e por meio desta bênção, Deus é
convidado a unir-se com os humanos e habitar a esfera do \emph{shabat}.
A noiva é o próprio \emph{shabat}, pois, como nos diz o
\emph{midrash}\footnote{O \emph{midrash} é uma metodologia judaica de
  ler as escrituras sagradas. Seu fundamento está em uma visão cíclica
  da história, em que o acontecimento central, ou melhor, a presença de
  Deus no texto, é continuamente revisado, oferecendo sempre novos
  significados às várias situações expostas, bem como às interpretações.}:
``Deus criou o tempo, bem como o dividiu em sete dias, e prometeu
ao \emph{shabat} (o sétimo dia da criação) que \emph{Israel será seu novo
par}\footnote{\emph{Midrash: Bereshit Rabá} 11.} \emph{para toda a
eternidade''.} Assim, todas as semanas o povo judeu saúda e celebra a
aproximação do \emph{shabat}, como se fosse um noivo que aguarda a
chegada de sua noiva.

Segundo a interpretação de Anat Yosef,\footnote{Anat Yosef foi um
  intérprete midráshico, isto é, comentarista das escrituras bíblicas.}
na verdade, a ``noiva'' mencionada pode ser entendida enquanto uma
alusão à própria \emph{Shechiná} --- a habitação e a presença de Deus.
Ou seja, durante o \emph{shabat} oramos pela aproximação da presença de
Deus no mundo humano. Por fim, segundo a filosofia judaica, se a
\emph{Shechiná} é a ``noiva'' --- que se aproxima do \emph{shabat} --- e
o povo judeu é o ``noivo'', então estamos a todo instante na esfera de
Deus, clamando pela presença de Sua habitação. Estamos em Deus, e
instamos para que a Sua existência esteja em nós também.

O Rabi Kobrin ensinava: ``Deus fala ao homem como falou a Moisés:
\emph{Tira os teus sapatos dos teus pés, descalça o habitual que encerra teus
pés e saberás que o lugar em que pisas é solo sagrado. Pois não há
situação da vida humana em que não se possa encontrar, em toda a parte e
a qualquer tempo, a santidade do Senhor}''\footnote{\textsc{buber}, \emph{op.\,cit.},
  p.\,490.}.

Por fim, se o universo --- e todas as coisas\footnote{Neste ensaio
  utiliza-se o sentido heideggeriano de ``coisas'', ou seja, as coisas
  como encontros de fios e de vida.} que ele contém, sejam elas vivas ou
mortas --- está dentro da esfera divina, então a esfera de Deus engloba
tanto o \emph{Olam Hazé} (``este mundo''), quanto o \emph{Olam Habá}
(``o mundo vindouro''), lembrando que a palavra hebraica \emph{Olam} é
usada tanto para designar ``tempo'', quanto para designar ``espaço''. E
Deus é o tempo e o espaço em que tudo acontece --- o lugar, a
completude, o eterno. Como nos lembra o pesquisador e estudioso das
escrituras bíblicas Jeff Benner:

\begin{quote}
Palavras hebraicas usadas para o espaço também são usadas para o tempo.
{[}\ldots{}{]} A palavra hebraica Olam significa literalmente ``além do
horizonte''. Ao olhar para longe, é difícil perceber detalhes e o que
está além desse horizonte não pode ser visto. Este conceito é o Olam. A
palavra Olam também é usada para o tempo: para o passado distante ou
para o futuro distante, como um tempo difícil de conhecer ou perceber.
Esta palavra é frequentemente traduzida como ``eternidade'', significa
um período de contínuo que nunca termina\footnote{BENNER,
  \emph{Eternity}.}.
\end{quote}

Assim, Deus não é infinito, mas eterno. Não há nem começo, nem fim em
Deus --- tudo é uma questão de meios (como são também as escrituras e as
interpretações do Talmud, ou as pedras dispostas sobre os
túmulos judaicos). O Deus judaico é a completude que habita os eventos e
os movimentos, e o judaísmo, então, é uma concepção do tempo.

\begin{center}\adforn{49}\end{center}

Voltemos ao Rabi Pinkhas: ``{[}\ldots{}{]} O grão semeado na terra
{[}\ldots{}{]} deve apodrecer, para que brote a nova espiga''\footnote{\textsc{buber},
  \emph{op.\,cit.}, p.166.}. Estes fragmentos de grãos, que são --- ao mesmo
tempo --- fragmentos de vida e de morte, são também sementes de
celebração. Estas sementes hão de morrer para renascer uma nova espiga.
E assim, Didi-Huberman nos lembra que estas sementes de renascimento de
vida devem ser guardadas ``não como tesouros imutáveis, mas como
sementes para o presente, para o futuro''\footnote{\textsc{didi-huberman}, op.
  cit., p.\,129.}.

Os poderes germinativos das sementes --- como aquelas encontradas nos
papéis do gueto de Varsóvia --- não foram armazenados para dar soluções
àqueles que as coletaram (essas pessoas já são cinzas\footnote{Estes
  indivíduos não sobreviveram, mas suas histórias sim.}). Ao contrário,
as sementes são respostas (não conclusivas) de vida para outrem, para o
presente e para o futuro --- para nós, e para a nossa geração. Por fim,
como nos lembra Rabi Pinkhas, se tudo nasce e tudo morre, para então
renascer, ou melhor, se ``{[}\ldots{}{]} No dia da destruição, o poder
jaz no fundo e aí cresce''\footnote{\textsc{buber}, \emph{op.\,cit.}, p.166.}
então, ``{[}\ldots{}{]} É por isso que nesse dia há de nascer o
Messias''\footnote{\emph{Ibidem}.}.

Entendendo que o ciclo nunca se completa sempre da mesma maneira, que
Deus habita o movimento\footnote{Deus habita os textos, os livros e o
  próprio trânsito da diáspora.} e ``os judeus, como tal, rejeitam
a estaticidade das coisas e das ideias, e acreditam na transformação e
na reedição''\footnote{\textsc{zevi}, \emph{Arquitetura e Judaísmo: Mendelsohn},
  p.\,10.}, nos permitimos aqui revisitar as nossas histórias --- de
judeus ou não judeus. Que possamos nos apropriar delas, da morte para o
renascimento, e da vida enquanto um princípio contínuo.

\begin{center}\adforn{49}\end{center}

\noindent{}``Quando o Rabi Abraão Iaakov, filho do Rabi Israel de Rijin,
ainda era criança, passeava ele no quintal, para cima e para baixo, ao
anoitecer de sexta-feira, quando os \emph{hassidim} já haviam ido rezar. Um
hassid dirigiu-se a ele e falou: \emph{Por que não entras? Já é \emph{shabat}}.
\emph{Ainda não é} --- respondeu Abraão. \emph{Como sabes isso?} --- perguntou o
hassid. \emph{No \emph{shabat}} --- replicou o menino --- \emph{surge sempre um novo céu
e ainda não o estou vendo}\footnote{\textsc{buber}, \emph{op.\,cit.}, p.\,383.}.

O calendário (e o seu tempo) é essencialmente lunar. Dessa forma, os
dias da semana, o \emph{shabat} e todos os feriados, as festas e as
festividades judaicas iniciam-se sempre no período de crepúsculo. Ao
final, o anoitecer permite um ``novo céu''\footnote{\emph{Ibidem}.}, um
novo recomeço. O dia judaico inicia ao nascer das três primeiras
estrelas no céu e termina no momento do pôr do sol.

As interpretações rabínicas sobre os textos sagrados não possuem
definições claras e concretas a respeito do momento do ``anoitecer''.
Porém, segundo o Talmud e o \emph{Midrash}, é nesse instante que
Deus, antes da entrega do \emph{shabat}, ao final do sexto dia da
criação divina, cria dez elementos, sendo eles: o arco-íris, o
\emph{maná}\footnote{\emph{Maná} é um alimento produzido milagrosamente,
  uma dádiva divina aos hebreus, durante toda a peregrinação no deserto
  do \emph{Sinai}, rumo à terra de \emph{Canaã}.}, a boca da terra, a
boca do jumento, as Tábuas da Lei, as letras, a escrita, o túmulo de
Moisés, o poço e o \emph{shamir}\footnote{Conta a tradição que o
  \emph{shamir} é um ser vivo criado por Deus e dotado da capacidade de
  cortar qualquer tipo de material, inclusive a mais dura rocha. O
  \emph{shamir} foi utilizado para a construção do Templo de Jerusalém,
  pois Deus não permitiria a utilização de arsenais de guerra para a
  construção deste lugar sagrado.}.

Entre o momento do pôr do sol e a aparição das três primeiras estrelas
na esfera celeste, os judeus se encontram em um limiar. Para a filosofia
judaica, nesse intervalo (e assim, nesse espaço-situação), não é nem dia
(em hebraico: \emph{iom}), nem noite (em hebraico: \emph{laila}), mas
sim um espaço-tempo entre o dia e a noite, chamado de \emph{erev} (em
português: ``tarde'').

\emph{Erev} é uma palavra hebraica cognata e da mesma raiz que as
palavras \emph{eruv} e \emph{arab}. \emph{Arab} significa ``mistura'',
\emph{eruv} é uma combinação de pátios e \emph{erev} é uma mescla entre
tempos, entre o dia e a noite. Esse limiar é conhecido como \emph{bein
hashemashot}, em português, ``entre sóis''.

Com uma duração de aproximadamente dezoito minutos (a depender da
localização geográfica), essa situação-período permite, por exemplo, o
acendimento das velas de \emph{shabat}\footnote{Por conta do judaísmo
  possuir um calendário lunar, todas as festas, festividades e
  celebrações judaicas iniciam-se no período do anoitecer, por meio do
  acendimento de fogo.}, em um espaço-tempo que não é ainda o
\emph{shabat}, e muito menos o sexto dia da criação divina, mas a
passagem entre estes dois períodos. Assim, \emph{bein hashemashot} é um
período, um espaço e um tempo limite e um limiar em si. Um momento entre
o dia e a noite, nem aqui, nem lá --- trânsito.

Como comenta Rabi Pinkhas: ``Despir uma forma, assumir uma forma,
isso se faz no instante do puro nada''\footnote{\textsc{buber}, \emph{op.\,cit.}, p.166.}.
A passagem entre o dia e a noite (a mudança e a transformação de uma
situação) se firma enquanto uma transição de forma, de matéria e de
espírito --- sair de um espaço e de um momento para entrar em outro.
Esta transição ou rito de passagem necessita de um tempo de pausa, de
reflexão, renovação e assim, de renascimento. Este movimento se faz no
``instante do puro nada''\footnote{\emph{Ibidem}.} --- no limiar e na
ambiguidade de conotações e de definições. O ``instante do puro
nada''\footnote{\emph{Ibidem}.} não é o vazio. O nada é a possibilidade.

\emph{Twilight People}, poema litúrgico do rabino norte-americano Reuven
Zellmann, se firma enquanto um lembrete de que este momento do
crepúsculo é um espaço plural e coletivo para todos. Tal espaço-tempo
sagrado de intermédio, instante de transformação e de renascimento,
reaparece em diversas religiões e crenças ao redor do mundo.

No poema do rabino Zellmann abençoa-se essa qualidade crepuscular dos
seres humanos. Abençoa-se a criação de um espaço-tempo de intermédio em
que as relações são interrompidas e o binarismo já não existe.
Abençoa-se o momento em que as crenças e as certezas absolutas são
suspensas e o julgamento é suavizado e esvaziado. Abençoa-se uma
situação de possibilidade, um momento do nada. E, finalmente, abençoa-se
Deus --- o criador do Universo, do dia e da noite, dos ventos, das
nuvens e do crepúsculo --- Aquele que transcende todas as categorias e
definições.

Entende-se aqui que a diáspora, ou os lugares judaicos, são espaços de
intermédio, que se firmam enquanto caminhos entrecruzados de memórias e
de destinos, de laços e de nós. A diáspora não pode ser determinada, mas
sempre pode ser vivida.

% \begin{center}\adforn{49}\end{center}